% Nouveau EGA V (cote 25): archive cote 25, site page 15 (PDF page 16).
% typescript. Draft compared: none.
% AI-transcribed by Claude Opus 5.5 (claude-opus-5-5, adaptive thinking, effort high), Message Batches API, batch ega25b (msgbatch_01A7n5YuXubUPNK1zsepi8RG), 2026-11-13; not yet checked by eye.
\scanpage{25/15}
\mspage{\hand{12}}

\uline{Corollaire} 2.17.   Soit Z l'ensemble des points de X en lesquels F
n'est pas séparable sur k (réf). Alors  $G_{\hand{\eta}}$ est séparable sur $k(\hand{\eta})$ sss
Z est fini, En particulier, si F est séparable sur k, $G_{\hand{\eta}}$ est séparable
sur $k(\hand{\eta})$.

\var{\margin{\unsure{bb pour|à mettre} dans l'énoncé 2.12.}}{\placedmargin{\unsure{bb pour|à mettre} dans l'énoncé 2.12.}}{marginal note placed}{ega25_25_015-v0001}\uline{Remarques} 2.18.  Le résultat clef  2.12. \var{\inserted{(et ses conséquences 2.13. à 2.17.)}}{(et ses conséquences 2.13. à 2.17.)}{insertion}{ega25_25_015-v0002} devient faux si on abandonne
l'hypothèse que f est une immersion, comme  on voit par exemple en
prenant pour  X un revêtement purement inséparable de P.  \var{\xout{Cxxx}}{}{typewriter strike-out}{ega25_25_015-v0003} Cepen-%
dant \var{\xout{lexrésultatxsuivantxx(quixauraitxdûxvenirxplûtôtxenxconclusion}}{}{typewriter strike-out}{ega25_25_015-v0004}
\var{\xout{àx2x11xx)xxxxxSixxX}}{}{typewriter strike-out}{ega25_25_015-v0005} \var{\xout{\unsure{Propxx}}}{}{typewriter strike-out}{ega25_25_015-v0006} si  k est \uline{de caractéristique nulle} ,  2.12. à
2.1\hand{7}. sont valables sans supposer f une immersion.  En effet, il suffit
de le vérifier pour 2.12. , et  cela résulte alors de 2.11. et du fait
que \var{\xout{xxxxx}}{}{typewriter strike-out}{ega25_25_015-v0007} pour un préschéma algébrique en c ar 0,  lisse = régulier.

\var{\deleted{Les propriétés envisagées jusqu'à maintenant étaient de nature
essentiellement locale. Nous allons terminer ce n° en étudiant les
propriétés globales}}{}{deletion}{ega25_25_015-v0008}

\var{\over{\unsure{3}}{20}}{20}{writing over}{ega25_25_015-v0009}.3.   \uline{Section hyperplane générique}: \uline{irréductibilité et connexité géomét.}
\uline{Théorème} 3.1. (Bertini-Zariski).  Supposons dim $f(X) \hand{\geq} 2$; et que X
soit \var{\xout{géométriquementxintègrex(resp.x}}{}{typewriter strike-out}{ega25_25_015-v0010} géométriquement irréductible\var{\deleted{)}}{}{deletion}{ega25_25_015-v0011}. Alors
la section hyperplane générique $Y_{\hand{\eta}}$ a la même propriété.

\var{\deleted{\var{\xout{x}}{}{typewriter strike-out}{ega25_25_015-v0013}(Compte tenu que X séparable \hspace{3em} Y séparable, les deux énoncés
de 3.1. sont en fait équivalents: \var{\xout{Onxpeutxdxailleursxsupposerx}}{}{typewriter strike-out}{ega25_25_015-v0014}  le
deuxième implique le premier, et inversement comme on voit en se rame-%
nant au cas k alg. clos.  Compte tenu de 2.1. et exmplicitant
la construction du corps des fonctions de Y \hspace{2em} comme extension de k( \hspace{1em} )
en termes de l'extension K de k et des éléments $t_1,\ldots,t_n$ de K images
d'un système de coordonnées affines convenable de}}{}{deletion}{ega25_25_015-v0012}

\mspage{\hand{15}}
